\documentclass[11pt]{article}
\usepackage[margin=1in]{geometry}
\usepackage{amsmath, amssymb, amsthm}
\usepackage{hyperref}
\usepackage{xcolor}

\newtheorem*{claim}{Claim}

% Visible placeholder for unfinished steps -- delete every \todo before submitting.
\newcommand{\todo}[1]{\textcolor{red}{\textbf{[TODO: #1]}}}
\newcommand{\blank}{\underline{\hspace{1.5cm}}}

\title{MAT 3150 --- Homework 4}
\author{}
\date{Due October 1}

\begin{document}
\maketitle

% ============================================================
\section*{Problem 1 (10 points)}
Let $(x_n)_n$ and $(y_n)_n$ be two sequences of real numbers, so that
$(y_n)_n$ is increasing. Assume that
\[
  \lim_{n\to\infty}\frac{x_{n+1}-x_n}{y_{n+1}-y_n} = L,
\]
for some $L\in\mathbb{R}$. Prove that for any $\varepsilon>0$ there exists
$N>0$ so that
\[
  (L-\varepsilon)(y_{n+1}-y_N) < x_{n+1}-x_N < (L+\varepsilon)(y_{n+1}-y_N),
\]
for all $n>N$.

\noindent\emph{Hint:} Use
$x_{n+1}-x_N = (x_{n+1}-x_n)+(x_n-x_{n-1})+\cdots+(x_{N+1}-x_N)$, for any $n>N$.

\begin{proof}
% Roadmap:
%  (a) Fix eps > 0. Write down what the hypothesis (the limit = L) gives you
%      for this eps -- an N and an inequality valid for k >= N.
%  (b) Unpack the absolute value into a two-sided inequality.
%  (c) Use the sign of y_{k+1} - y_k (what does "increasing" give you?) to
%      clear the denominator. Be careful about which way the inequality goes.
%  (d) Apply (c) to each k = N, N+1, ..., n and combine using the hint.
%  (e) Identify what the y-terms collapse to.

% Reading the goal: for n > N, the quantity x_{n+1} - x_N is trapped between
%   lower bound (L - eps)(y_{n+1} - y_N)  and  upper bound (L + eps)(y_{n+1} - y_N).
% The hypothesis only bounds ONE step (the ratio for a single k), so the plan is:
%   one-step bound on the ratio -> one-step bound on x_{k+1} - x_k -> sum the steps.

Let $\varepsilon>0$. For each $k$, write
\[
  r_k = \frac{x_{k+1}-x_k}{y_{k+1}-y_k},
\]
so the hypothesis says $\lim_{k\to\infty} r_k = L$.

\textbf{Step 1: Use the hypothesis.}
By the definition of $\lim_{k\to\infty} r_k = L$, there exists $N$ such that
for all $k \ge N$,
\[
  |r_k - L| < \varepsilon .
\]

Unpacking the absolute value, $-\varepsilon < r_k - L < \varepsilon$, and
adding $L$ throughout gives, for all $k\ge N$,
\[
  L - \varepsilon < \frac{x_{k+1}-x_k}{y_{k+1}-y_k} < L + \varepsilon .
\]

\textbf{Step 2: Clear the denominator.}
We multiply the inequality from Step 1 by $y_{k+1}-y_k$. Since
$y_{k+1}-y_k > 0$, multiplying by it preserves the direction of the strict
inequalities $<$.

Indeed, since $(y_n)$ is increasing, $y_{k+1}-y_k \ge 0$; and
$y_{k+1}-y_k \ne 0$, since otherwise the fraction in the hypothesis would be
undefined. Hence $y_{k+1}-y_k>0$, and we obtain, for all $k\ge N$,
\[
  (L-\varepsilon)(y_{k+1}-y_k) < x_{k+1}-x_k < (L+\varepsilon)(y_{k+1}-y_k).
\]

\textbf{Step 3: Sum from $k=N$ to $k=n$.}
Let $n>N$. The inequality from Step 2 holds for each $k=N,N+1,\dots,n$.
Adding these $n-N+1$ inequalities gives
\[
  \sum_{k=N}^{n}(L-\varepsilon)(y_{k+1}-y_k)
  < \sum_{k=N}^{n}(x_{k+1}-x_k)
  < \sum_{k=N}^{n}(L+\varepsilon)(y_{k+1}-y_k).
\]
Since $L-\varepsilon$ and $L+\varepsilon$ do not depend on $k$, we can pull
them out of the sums:
\[
  (L-\varepsilon)\sum_{k=N}^{n}(y_{k+1}-y_k)
  < \sum_{k=N}^{n}(x_{k+1}-x_k)
  < (L+\varepsilon)\sum_{k=N}^{n}(y_{k+1}-y_k).
\]

\textbf{Step 4: Simplify.}
The $y$-sum telescopes, since every $y_{N+1},\dots,y_n$ appears once with $+$
and once with $-$:
\[
  \sum_{k=N}^{n}(y_{k+1}-y_k) = y_{n+1}-y_N .
\]
Likewise, the $x$-sum telescopes (this is the hint):
\[
  \sum_{k=N}^{n}(x_{k+1}-x_k) = x_{n+1}-x_N .
\]

Substituting both sums into the inequality from Step 3 gives
\[
  (L-\varepsilon)(y_{n+1}-y_N) < x_{n+1}-x_N < (L+\varepsilon)(y_{n+1}-y_N)
\]
for all $n>N$, which is what we wanted to prove.

\end{proof}

% ============================================================
\section*{Problem 2 (10 points)}
Let $(x_n)_n$ and $(y_n)_n$ be two sequences of real numbers, so that
$(y_n)_n$ is increasing and $\lim_{n\to\infty}y_n=\infty$. Assume that
\[
  \lim_{n\to\infty}\frac{x_{n+1}-x_n}{y_{n+1}-y_n} = L,
\]
for some $L\in\mathbb{R}$. Prove that
\[
  \lim_{n\to\infty}\frac{x_n}{y_n}
\]
exists and is also equal to $L$.

\noindent\emph{Hint:} Use Problem 1.

\begin{proof}
% Roadmap:
%  (a) Goal: for every eps > 0, find M with |x_n/y_n - L| < (some multiple of) eps
%      for n >= M.
%  (b) Start from Problem 1's inequality. You want x_{n+1}/y_{n+1} in the
%      middle -- what do you divide by, and why is that allowed for large n?
%  (c) After dividing, some terms involve the FIXED index N. What happens to
%      them as n -> infinity (use y_n -> infinity)?
%  (d) Turn "those terms go to 0" into a concrete second threshold, and take
%      the larger of the two thresholds.

Let $\varepsilon>0$, and let $N$ be as in Problem 1.

\textbf{Step 1: Rearrange Problem 1's inequality.}
By Problem 1, for all $n>N$,
\[
  (L-\varepsilon)(y_{n+1}-y_N) < x_{n+1}-x_N < (L+\varepsilon)(y_{n+1}-y_N).
\]
Adding $x_N$ throughout,
\[
  x_N+(L-\varepsilon)(y_{n+1}-y_N) < x_{n+1} < x_N+(L+\varepsilon)(y_{n+1}-y_N).
\]
Expanding the left half,
\[
  x_N+(L-\varepsilon)y_{n+1}-(L-\varepsilon)y_N < x_{n+1} .
\]
and the right half,
\[
  x_{n+1} < x_N+(L+\varepsilon)y_{n+1}-(L+\varepsilon)y_N .
\]

\todo{Why is $y_{n+1}>0$ for large $n$? Use $y_n\to\infty$ to get an index $N_0$ with $y_{n+1}>0$ for all $n\ge N_0$.}

Dividing the left half by $y_{n+1}>0$,
\[
  L-\varepsilon + \frac{x_N-(L-\varepsilon)y_N}{y_{n+1}} < \frac{x_{n+1}}{y_{n+1}} .
\]
and the right half,
\[
  \frac{x_{n+1}}{y_{n+1}} < L+\varepsilon + \frac{x_N-(L+\varepsilon)y_N}{y_{n+1}} .
\]
Together,
\[
  L-\varepsilon + \frac{x_N-(L-\varepsilon)y_N}{y_{n+1}}
  < \frac{x_{n+1}}{y_{n+1}}
  < L+\varepsilon + \frac{x_N-(L+\varepsilon)y_N}{y_{n+1}} .
\]

\textbf{Step 2: Handle the terms involving $N$.}
The leftover terms are $\dfrac{A}{y_{n+1}}$ and $\dfrac{B}{y_{n+1}}$, where
\[
  A = x_N-(L-\varepsilon)y_N, \qquad B = x_N-(L+\varepsilon)y_N .
\]
Since $N$ is fixed, $A$ and $B$ are constants: they do not depend on $n$.
Since $y_{n+1}\to\infty$, a constant divided by $y_{n+1}$ tends to $0$, so
\[
  \frac{A}{y_{n+1}}\to 0 \qquad\text{and}\qquad \frac{B}{y_{n+1}}\to 0 .
\]

\textbf{Step 3: Choose the final threshold.}
Since $\frac{A}{y_{n+1}}\to0$, there exists $M_1$ such that for $n\ge M_1$,
\[
  \left|\frac{x_N-(L-\varepsilon)y_N}{y_{n+1}}\right| < \varepsilon ,
\]
and since $\frac{B}{y_{n+1}}\to0$, there exists $M_2$ such that for $n\ge M_2$,
\[
  \left|\frac{x_N-(L+\varepsilon)y_N}{y_{n+1}}\right| < \varepsilon .
\]
Let $M=\max\{N, N_0, M_1, M_2\}$, where $N_0$ is the index from Step 1 with
$y_{n+1}>0$ for $n\ge N_0$.

\todo{For $n\ge M$, plug these bounds into Step 1: \ \blank\ $<\frac{x_{n+1}}{y_{n+1}}<$ \blank.}

\todo{Conclude (Fix 1: the bound is $2\varepsilon$ and $\varepsilon$ was arbitrary; or Fix 2: redo with $\varepsilon/2$ from the start).}

\end{proof}

% ============================================================
\section*{Problem 3 (10 points)}
Prove that if
\[
  \lim_{n\to\infty}a_n = L,
\]
then
\[
  \lim_{n\to\infty}\frac{a_1+a_2+\cdots+a_n}{n}
\]
exists and equals $L$.

\begin{proof}
% Roadmap:
%  Question to ask yourself: the expression is a quotient x_n / y_n.
%  What would x_n and y_n be? Do they satisfy the hypotheses of Problem 2?
%  What is (x_{n+1}-x_n)/(y_{n+1}-y_n)?

\todo{Choose $x_n = \blank$ and $y_n = \blank$.}

\todo{Check the hypotheses of Problem 2: $(y_n)$ increasing, $y_n\to\infty$.}

\todo{Compute $\dfrac{x_{n+1}-x_n}{y_{n+1}-y_n}$ and find its limit.}

\todo{Conclude.}

\end{proof}

% ============================================================
\section*{Problem 4 (10 points)}
Define the sequence $a_1=1$ and
\[
  a_{n+1} = 3-\frac{1}{a_n}.
\]
Show that $a_n$ is convergent, by proving that $a_n$ is monotone and bounded.
What is the limit of $a_n$?

% Scratch work (not part of the write-up): compute a_2, a_3, a_4 to guess
% whether the sequence is increasing or decreasing, and guess a bound.

\begin{claim}
$(a_n)$ is \blank\ (increasing / decreasing) and $\blank \le a_n \le \blank$ for all $n$.
\end{claim}

\begin{proof}
% Roadmap:
%  (a) Bounded: induction. Base case n = 1; inductive step a_n -> a_{n+1}.
%      Note you need a_n > 0 for 1/a_n to behave nicely -- make sure your
%      bound guarantees that.
%  (b) Monotone: induction again, OR compute a_{n+1} - a_n directly.
%  (c) Apply the Monotone Convergence Theorem.
%  (d) Find the limit: let a = lim a_n and pass to the limit in the recursion.
%      You'll get more than one candidate -- use the bounds to pick the right one.

\textbf{Step 1: Bounded.}
\todo{Base case.}
\todo{Inductive step: assume $\blank \le a_n \le \blank$, show the same for $a_{n+1}$.}

\textbf{Step 2: Monotone.}
\todo{Base case: compare $a_1$ and $a_2$.}
\todo{Inductive step.}

\textbf{Step 3: Convergence.}
\todo{Cite the Monotone Convergence Theorem.}

\textbf{Step 4: The limit.}
\todo{Let $a=\lim a_n$. Why is $a \neq 0$? Pass to the limit in $a_{n+1}=3-\frac1{a_n}$ and solve.}
\todo{Which root, and why?}

\end{proof}

\end{document}
